% Theta Complex-Time Bridge
% Status: Bridge / non-canonical derivation
% Purpose: rigorous bridge note connecting a reduced Jacobi-theta complex-time
% construction to the existing UBT complex-time and operator language.

\documentclass[11pt,a4paper]{article}
\usepackage{amsmath,amssymb,amsthm}
\usepackage{hyperref}
\usepackage{geometry}
\geometry{margin=2.5cm}

\begin{document}

\section{Theta Transform with Complex Time as a Bridge to UBT}
\label{sec:theta_complex_time_bridge}

\paragraph{Status.}
Bridge / non-canonical derivation. This note is a mathematically motivated bridge between a reduced Jacobi-theta construction and the canonical UBT framework. It does not replace the canonical definitions of $T_B$ or $\Theta(q,T_B)$.

\subsection{Motivation}

\paragraph{Canonical anchor.}
Canonical UBT uses biquaternion time
\begin{equation}
T_B = t + i\psi + j\chi + k\xi,
\end{equation}
with isotropic complex-time limit
\begin{equation}
\tau = t + i\psi.
\end{equation}

\paragraph{Bridge statement.}
The Jacobi-theta complex-time kernel provides a reduced model in which imaginary time gives Gaussian damping and real time gives quadratic phase evolution. This structure is compatible with heat-kernel, diffusion, and spectral viewpoints already present in UBT.

\paragraph{Conjectural boundary.}
This note does not derive the full UBT field, metric sector, or Standard Model sector from the theta construction.

\subsection{Jacobi Theta Starting Point}

\paragraph{Bridge statement.}
We begin from the classical Jacobi theta function
\begin{equation}
\theta(z,\tau)
=
\sum_{n \in \mathbb{Z}}
\exp\!\left(
\pi i \tau n^2 + 2\pi i n z
\right),
\qquad \Im(\tau) > 0.
\end{equation}

Introduce the reduced complex-time parameter
\begin{equation}
\tau = t + i\phi.
\end{equation}

Then the quadratic exponent decomposes as
\begin{equation}
\exp(\pi i \tau n^2)
=
\exp(-\pi \phi n^2)\,\exp(\pi i t n^2).
\end{equation}

\paragraph{Interpretation.}
The factor $\exp(-\pi \phi n^2)$ is heat-kernel-like damping, while $\exp(\pi i t n^2)$ is oscillatory quadratic phase evolution.

\subsection{Reduced Discrete Theta Projection}

\paragraph{Bridge statement.}
For a reduced discrete mode sequence $a_n^{(s)}$, define
\begin{equation}
S_s(t,\phi)
=
\sum_{n=0}^{N-1}
a_n^{(s)}\,\exp(-\pi \phi n^2)\,\exp(\pi i t n^2).
\end{equation}

\paragraph{Interpretation.}
The symbol $s$ is an external evolution label for the reduced family of mode sequences. It is distinct from the real kernel time $t$ and distinct from the imaginary damping scale $\phi$.

\subsection{Energy Observable}

\paragraph{Bridge statement.}
Define the energy-like observable
\begin{equation}
\Theta_s(t,\phi) = \left|S_s(t,\phi)\right|^2.
\end{equation}

Expanding the modulus square gives
\begin{equation}
\Theta_s(t,\phi)
=
\sum_{n,m}
a_n^{(s)}\overline{a_m^{(s)}}\,
\exp\!\left(-\pi \phi (n^2+m^2)\right)\,
\exp\!\left(\pi i t (n^2-m^2)\right).
\end{equation}

\paragraph{Interpretation.}
The double sum makes pairwise mode interference explicit. The single-sum amplitude and double-sum energy forms are equivalent descriptions of the same reduced observable sector.

\subsection{Averaging over Imaginary Time}

\paragraph{Bridge statement.}
To reduce sensitivity to any single damping scale, define the averaged observable
\begin{equation}
\overline{\Theta}_s(t)
=
\frac{1}{P}\sum_{j=1}^{P}\Theta_s(t,\phi_j).
\end{equation}

\paragraph{Interpretation.}
This averaging plays the role of a scale-stabilized reduced observable over the real kernel time $t$.

\subsection{Alignment with Canonical UBT Notation}

\paragraph{Canonical anchor.}
In canonical UBT notation, the scalar imaginary component of complex time is $\psi$, not $\phi$, and the isotropic limit is
\begin{equation}
\tau = t + i\psi.
\end{equation}

\paragraph{Bridge statement.}
The reduced theta notation therefore maps
\begin{equation}
\phi \longleftrightarrow \psi
\end{equation}
only as a notation bridge inside the isotropic complex-time limit.

The reduced amplitude may then be rewritten in UBT-aligned notation as
\begin{equation}
S_s(t,\psi)
=
\sum_{n=0}^{N-1}
a_n^{(s)} e^{-\pi \psi n^2} e^{\pi i t n^2}.
\end{equation}

\paragraph{Interpretation.}
This notation change is a compatibility step with canonical UBT. It is not a new canonical definition and does not extend the reduced model to full biquaternion time.

\paragraph{Canonical anchor.}
The symbol $\Theta$ is already reserved in UBT for the fundamental field $\Theta(q,T_B)$.

\paragraph{Bridge statement.}
To avoid collision with the canonical field symbol, it is preferable in UBT-facing writing to rename the reduced observable
\begin{equation}
\Theta_s(t,\psi)
\end{equation}
as
\begin{equation}
\mathcal{E}_s(t,\psi) := \left|S_s(t,\psi)\right|^2,
\qquad
\overline{\mathcal{E}}_s(t) := \frac{1}{P}\sum_{j=1}^{P}\mathcal{E}_s(t,\psi_j).
\end{equation}

\paragraph{Interpretation.}
This is a notation hygiene step only. The reduced mathematics is unchanged.

\subsection{Operator and Spectral Viewpoint}

\paragraph{Canonical anchor.}
UBT already uses operator language through covariant derivatives, Laplace-type structures, and spectral arguments in the complex-time sector.

\paragraph{Bridge statement.}
The kernel
\begin{equation}
e^{-\pi \psi n^2} e^{\pi i t n^2}
\end{equation}
has the form of a discrete complex-time propagator factor: Gaussian damping in imaginary time and oscillatory phase evolution in real time.

\paragraph{Interpretation.}
This suggests that $S_s(t,\psi)$ may be read as a reduced projected propagator-like amplitude on a compact or discrete mode space, while $\overline{\mathcal{E}}_s(t)$ is a reduced observable derived from it.

\paragraph{Conjectural boundary.}
The identification of this reduced amplitude with the full UBT field propagator is not proved here.

\subsection{Limits of the Bridge}

\paragraph{Canonical anchor.}
The canonical field remains $\Theta(q,T_B)$ and the canonical time remains $T_B = t + i\psi + j\chi + k\xi$.

\paragraph{Bridge statement.}
The theta construction presented here is a reduced model compatible with the isotropic complex-time limit of UBT.

\paragraph{Conjectural boundary.}
This note does not prove exact recovery of the metric sector, the GR limit, the Standard Model interaction sectors, or the full biquaternion-valued dynamics from the reduced theta model.

\subsection{Next Steps}

\paragraph{Bridge step.}
Derive an operator form whose matrix elements reproduce $S_s(t,\psi)$ as a reduced amplitude.

\paragraph{Bridge step.}
Test whether $\overline{\mathcal{E}}_s(t)$ behaves like a reduced propagator observable in mode spaces already natural in UBT.

\paragraph{Bridge step.}
Compare the reduced theta kernel with existing UBT complex-time and diffusion-like structures.

\paragraph{Conjectural step.}
Determine whether the quadratic phase factor can be embedded into full biquaternion-valued evolution.

\end{document}